\begin{frame}[plain]
  \centering
  \begin{columns}[c]
    \begin{column}{0.34\textwidth}
      \centering
      \includegraphics[width=0.72\linewidth]{fisher_portrait.jpg}\\[0.3em]
      {\small 로널드 피셔 (1952) \\[0.15em] 통계학 컨퍼런스, Asheville NC}
    \end{column}
    \begin{column}{0.32\textwidth}
      \centering
      \includegraphics[width=0.95\linewidth]{iris_setosa.jpg}\\[0.3em]
      {\small \textit{Iris setosa}}
    \end{column}
    \begin{column}{0.32\textwidth}
      \centering
      \includegraphics[width=0.95\linewidth]{iris_versicolor.jpg}\\[0.3em]
      {\small \textit{Iris versicolor}}
    \end{column}
  \end{columns}
  \vskip0.4em
  {\footnotesize 1936년 ``분류 문제에서 다중 측정의 이용'' --- 붓꽃(iris) 종을
  꽃술·꽃잎 측정값으로 분류한 데이터셋이 오늘날 \kb{피셔 선형 판별}의 기원}
\end{frame}
